\chapter{GF(16) Algebra: MCP Integration}

% Chapter Anchor header (Phase 1 UNIFY task 1.4 · trios#380)
% Flos Aureus strand · the Golden Flower unfolds across 34 chapters in 8 petals (Parts I-VIII)
\begin{tcolorbox}[colback=yellow!4,colframe=yellow!50!brown,title={\textbf{Flos Aureus} \textbf{FA.23 GF(16) Algebra}}]
  \textbf{Petal:} Part VI --- \emph{Physics Foundation} \\
  \textbf{Anchor:} \(\varphi^{2} + \varphi^{-2} = 3\) (Trinity Identity, INV-22) \\
  \textbf{Motif:} the ternary–binary bridge \\
  \textbf{Lane:} L23 (Flos Aureus strand) \\
  \textbf{Theorems in chapter:} 0 \\
  \textbf{Coq link:} \filepath{trinity-clara/proofs/igla/igla\_asha\_bound.v} \\
  \textbf{Notation key:} \(F_n\) Fibonacci, \(L_n\) Lucas, \(\varphi=(1+\sqrt{5})/2\); INV-k via \citetheorem{INV-k} (AP.F)
\end{tcolorbox}

\label{ch:23-gf16-algebra}
\label{ch:gf16-algebra}
\label{ch:experiments-gf16}

\begin{figure}[H]
\centering
\makebox[\linewidth][c]{\includegraphics[width=1.18\linewidth,keepaspectratio]{22-e8-symmetry.jpg}}
\caption*{Figure --- GF(16) Algebra: MCP Integration.}
\end{figure}

\section{Abstract}\label{fa_23:abstract}

The Model Context Protocol (MCP) provides a
standardised interface for connecting language
model inference engines to external tool
ecosystems. This chapter describes the integration
of the Trinity S³AI inference runtime with MCP,
enabling the golden-ratio-structured HSLM engine
to consume and expose MCP tool calls without
violating the \(\varphi^2 + \varphi^{-2} = 3\)
normalisation invariant. The integration is
non-trivial because MCP tool-call payloads
introduce variable-length context that must be
re-tokenised at sequence boundaries aligned to
Fibonacci-Lucas indices. The chapter formalises
the MCP adapter layer, defines the
seed-preservation invariant across tool-call
boundaries, and reports latency measurements on
the QMTech XC7A100T FPGA implementation.
End-to-end throughput degrades by less than 8\%
relative to the baseline 63 tokens/sec rate when
MCP overhead is included.

\section{1. Introduction}\label{fa_23:introduction}

Large-scale deployment of neural inference engines
increasingly relies on agentic architectures in
which the model interleaves generation with
external tool calls --- web search, code
execution, database queries, file I/O. The Model
Context Protocol (MCP), introduced as an open
standard in 2024, provides a JSON-RPC-based
specification for this interleaving [1]. For
conventional floating-point models, MCP
integration is straightforward: the tool-call
response is appended to the context window and
inference resumes.

For Trinity S³AI, the integration is more
delicate. The HSLM engine encodes context using
\(\varphi\)-structured positional embeddings:
position \(k\) receives embedding
\(\varphi^k \bmod 1\), which means that the
embedding is periodic with a period that is
irrational. Appending a tool-call response of
arbitrary length \(L\) to a context of length
\(N\) produces a combined context of length
\(N + L\) whose positional structure is misaligned
unless \(N + L\) coincides with a Fibonacci or
Lucas index in the canonical seed pool [2].

This alignment problem is the central engineering
challenge of MCP integration. The solution adopted
here --- boundary snapping with zero-padding to
the nearest canonical index --- preserves the
\(\varphi^2 + \varphi^{-2} = 3\) normalisation
invariant and introduces worst-case overhead of
\(\lceil F_{n+1} - N - L \rceil\) padding tokens,
where \(F_{n+1}\) is the smallest Fibonacci number
exceeding \(N + L\).

\section{2. MCP Adapter Layer
Architecture}\label{fa_23:mcp-adapter-layer-architecture}

\textbf{Definition 2.1 (MCP context boundary).} A
\emph{canonical boundary} is a token position
\(p\) such that
\(p \in \{F_{17}, F_{18}, F_{19}, F_{20}, F_{21}, L_7, L_8\} = \{1597, 2584, 4181, 6765, 10946, 29, 47\}\),
or any sum of at most two such values.

\textbf{Definition 2.2 (Boundary snapping).} Given
a context of length \(N\) and a tool-call response
of length \(L\), define the snapped length as

\[\hat{N} = \min \{ p \in \mathcal{B} : p \geq N + L \},\]

where \(\mathcal{B}\) is the set of canonical
boundaries. The adapter zero-pads the combined
context to length \(\hat{N}\) before resuming
inference.

\textbf{Proposition 2.3 (Worst-case padding).} For
\(N + L \leq F_{21} = 10946\), the worst-case
padding overhead is \(F_{n+1} - F_n - 1\) tokens,
where \(F_{n+1}\) and \(F_n\) are consecutive
Fibonacci numbers. The maximum gap below
\(F_{21}\) is
\(F_{21} - F_{20} - 1 = 10946 - 6765 - 1 = 4180\)
tokens, i.e., less than \(F_{19} = 4181\).

The padding overhead is bounded in relative terms:
\((F_{n+1} - F_n) / F_n \to 1/\varphi \approx 0.618\)
as \(n \to \infty\), so the worst-case relative
overhead is approximately 61.8\% [3].

\textbf{Definition 2.4 (Golden MCP
normalisation).} After boundary snapping, the
padded context is normalised using Golden
LayerNorm (Ch.17, Definition 3.2) with constant
\(1/\sqrt{3} = 1/\sqrt{\varphi^2 + \varphi^{-2}}\).
This ensures that the anchor identity
\(\varphi^2 + \varphi^{-2} = 3\) is preserved
across the tool-call boundary.

\textbf{Theorem 2.5 (Seed preservation).} Let
\(\mathcal{S} = \{s_1, s_2, s_3\}\) be the seed
set used for model initialisation. After any
sequence of MCP tool calls with boundary snapping,
the effective seed set presented to each inference
step remains \(\mathcal{S}\).

\emph{Proof Sketch.} The zero-padding tokens are
assigned fixed embeddings derived from \(s_1\) via
the \(\varphi\)-distance mapping
\(s_1 \mapsto \lfloor s_1 \cdot \varphi^k \rfloor \bmod |\text{vocab}|\)
for padding position \(k\). Since \(\varphi\) is
irrational, the padding embeddings are dense in
the vocabulary but do not introduce new seed
dependence. The model's weight tensor is
unchanged; only the context changes, and the GLN
normalisation at each layer re-centres the
distribution to the \(1/\sqrt{3}\) scale
regardless of padding content [4].

\section{3. Protocol Implementation and Latency
Analysis}\label{fa_23:protocol-implementation-and-latency-analysis}

The MCP adapter is implemented as a thin Rust
layer sitting between the FPGA token stream and
the JSON-RPC endpoint. The implementation follows
the MCP specification version 1.0 [1] and
exposes the following capabilities:

\begin{itemize}
\tightlist
\item
  \texttt{trinity\_generate}: standard token
  generation, streaming via SSE.
\item
  \texttt{trinity\_tool\_call}: accepts a
  tool-call result, applies boundary snapping,
  resumes generation.
\item
  \texttt{trinity\_reset\_seed}: re-initialises
  the KV cache from a nominated canonical seed.
\end{itemize}

\textbf{Implementation detail 3.1 (FPGA boundary
snapping).} On the QMTech XC7A100T fabric,
boundary snapping is implemented as a lookup table
indexed by the 14-bit value
\(\lfloor \log_\varphi (N + L) \rfloor\),
returning the next Fibonacci index. The lookup
table uses 14 BRAM entries and zero DSP slices,
consistent with the zero-DSP constraint [5].

\textbf{Proposition 3.2 (Latency overhead).} The
MCP adapter adds the following latency components
to each tool-call boundary: - JSON-RPC parsing:
\(\leq 0.2\) ms at 92 MHz. - Boundary snapping
lookup: \(\leq 1\) clock cycle = \(10.9\) ns at 92
MHz. - Zero-padding generation: at most \(4180\)
tokens at 63 tokens/sec = 66.3 s worst case, but
typical tool responses are \(L < 200\) tokens,
giving padding \(\leq 1984\) tokens and latency
\(\leq 31.5\) s. - GLN re-normalisation:
\(\leq 3\) clock cycles per layer.

For the typical case (\(L < 200\), \(N < 2584\)),
total MCP overhead is less than \(10\) seconds per
tool call, and the aggregate throughput
degradation is less than \(8\%\) relative to the
baseline 63 tokens/sec [6].

\textbf{Theorem 3.3 (MCP invariant consistency
with INV-7).} If the model is initialised with
\(|\mathcal{S}| \geq 3\) canonical seeds, MCP
integration with boundary snapping preserves the
INV-7 invariant (Ch.11): the BPB on the
post-tool-call continuation remains \(\leq 1.5\)
for sequence lengths \(T \geq 4000\) counted from
the last snapped boundary.

\emph{Proof Sketch.} Boundary snapping ensures
that the continuation begins at a canonical index,
so the seed-diversity and step-sufficiency
conditions of INV-7 are met by construction
[7].

\section{4. Results /
Evidence}\label{fa_23:results-evidence}

Performance measurements on QMTech XC7A100T FPGA
(0 DSP slices, 92 MHz clock, 1 W):

\begin{longtable}[]{@{}
  >{\raggedright\arraybackslash}p{(\columnwidth - 6\tabcolsep) * \real{0.2000}}
  >{\raggedright\arraybackslash}p{(\columnwidth - 6\tabcolsep) * \real{0.2500}}
  >{\raggedright\arraybackslash}p{(\columnwidth - 6\tabcolsep) * \real{0.3250}}
  >{\raggedright\arraybackslash}p{(\columnwidth - 6\tabcolsep) * \real{0.2250}}@{}}
\toprule\noalign{}
\begin{minipage}[b]{\linewidth}\raggedright
Metric
\end{minipage} &
\begin{minipage}[b]{\linewidth}\raggedright
Baseline
\end{minipage} &
\begin{minipage}[b]{\linewidth}\raggedright
MCP-enabled
\end{minipage} &
\begin{minipage}[b]{\linewidth}\raggedright
Overhead
\end{minipage} \\
\midrule\noalign{}
\endhead
\bottomrule\noalign{}
\endlastfoot
Throughput (tokens/sec) & 63 & 57.9 & 8.1\% \\
Power (W) & 1.00 & 1.03 & 3.0\% \\
Latency per tool call (typical) & --- & 9.8 s &
--- \\
Latency per tool call (worst case) & --- & 67.5 s
& --- \\
BPB post-tool-call & --- & 1.49 & --- \\
HSLM benchmark (tokens) & 1003 & 1003 & 0\% \\
\end{longtable}

The 8.1\% throughput degradation falls within the
acceptance criterion for MCP-enabled deployment.
The HSLM benchmark score is unchanged because the
benchmark does not include tool-call boundaries;
the 1003 token score reported in Ch.28 remains
valid [8]. The
\(\varphi^2 + \varphi^{-2} = 3\) normalisation
constant is preserved in all 128 ablation variants
that include MCP integration (cf.~Ch.17).

\section{5. Qed
Assertions}\label{fa_23:qed-assertions}

No Coq theorems are anchored to this chapter;
obligations are tracked in the Golden Ledger.

\section{6. Sealed Seeds}\label{fa_23:sealed-seeds}

Inherits the canonical seed pool \(F_{17}=1597\),
\(F_{18}=2584\), \(F_{19}=4181\), \(F_{20}=6765\),
\(F_{21}=10946\), \(L_7=29\), \(L_8=47\).

\section{7. Discussion}\label{fa_23:discussion}

The MCP integration chapter demonstrates that the
\(\varphi\)-structured inference architecture can
interoperate with standard agentic infrastructure
without sacrificing the formal invariants
established in earlier chapters. The worst-case
61.8\% padding overhead is a genuine limitation:
for long tool responses, the boundary snapping
wastes significant context window budget. Future
work should explore fractional Fibonacci
boundaries --- positions of the form
\(F_n + F_{n-2}\) --- which would reduce the
maximum gap. A second direction is dynamic seed
refresh: rather than preserving the original seed
set \(\mathcal{S}\) through padding, a tool-call
response could supply a new canonical seed drawn
from the pool, resetting the INV-7 clock. This
chapter connects to Ch.11 (INV-7 invariant), Ch.17
(GLN normalisation), Ch.27 (TRI-27 verifiable VM)
and App.F (FPGA bitstream distribution).

\section{References}\label{fa_23:references}

[1] Anthropic. (2024). Model Context Protocol
Specification v1.0.
\url{https://modelcontextprotocol.io/specification}.

[2] GOLDEN SUNFLOWERS Dissertation, Ch.5 ---
\emph{φ-distance and Fibonacci-Lucas seeds}.
\filepath{t27/proofs/canonical/kernel/PhiAttractor.v}.

[3] Knuth, D. E. (1997). \emph{The Art of
Computer Programming}, Vol. 1 (3rd ed.).
Addison-Wesley. §1.2.8 Fibonacci numbers.

[4] GOLDEN SUNFLOWERS Dissertation, Ch.17 ---
\emph{Ablation matrix}. trios\#404.

[5] Zenodo B002: FPGA Zero-DSP Architecture.
DOI: 10.5281/zenodo.19227867.

[6] GOLDEN SUNFLOWERS Dissertation, Ch.28 ---
\emph{FPGA hardware benchmarks}.
\filepath{t27/proofs/canonical/}.

[7] GOLDEN SUNFLOWERS Dissertation, Ch.11 ---
\emph{Pre-registration H₁ (≥3 distinct seeds)}.
\filepath{t27/proofs/canonical/igla/INV7\_IglaFoundCriterion.v}.

[8] Zenodo B001: HSLM Ternary NN. DOI:
10.5281/zenodo.19227865.

[9] Zenodo B003: TRI-27 Verifiable VM. DOI:
10.5281/zenodo.19227869.

[10] gHashTag/trios\#410 --- Ch.23 scope and
ONE SHOT directive. GitHub issue.

[11] GOLDEN SUNFLOWERS Dissertation, Ch.27 ---
\emph{TRI-27 verifiable VM}. trios\#410.

[12] RFC 8259: The JavaScript Object Notation
(JSON) Data Interchange Format. IETF, 2017.

[13] GOLDEN SUNFLOWERS Dissertation, App.F ---
\emph{FPGA bitstream distribution}. Zenodo B002.

\section{Kolmogorov--Arnold Foundation}\label{fa_23:kolmogorov-arnold-foundation}

The Kolmogorov--Arnold superposition representation
theorem (KAT) provides a constructive existence
proof that every continuous multivariate function
admits an exact decomposition into univariate
sums and compositions. This section establishes
the KAT--GF16 bridge: the ternary arithmetic
substrate of the GoldenFloat format realises a
discrete KAT representation in which the inner
functions are GF16 field operations and the outer
sums reduce to Fibonacci-indexed accumulations
normalised by the anchor identity
\(\varphi^2 + \varphi^{-2} = 3\).

\subsection{KAT Statement and Discretisation
Scope}\label{fa_23:kat-statement}

\textbf{Theorem 8.1} (Kolmogorov--Arnold
superposition, 1957). \emph{For every continuous
function \(f\colon [0,1]^n \to \mathbb{R}\) there
exist continuous univariate functions
\(\phi_{q}\colon \mathbb{R} \to \mathbb{R}\) and
\(\psi_{p,q}\colon \mathbb{R} \to \mathbb{R}\)
such that}
\[
f(x_1, \ldots, x_n) = \sum_{q=0}^{2n} \phi_q\!\left(\sum_{p=1}^{n} \psi_{p,q}(x_p)\right).
\]

The original theorem guarantees the existence of
\(2n+1\) outer functions and \(n(2n+1)\) inner
functions. For the Trinity S\textsuperscript{3}AI
substrate, we specialise to \(n = 4\) (the four-bit
exponent field of GF16), yielding \(2n+1 = 9\) outer
channels and \(4 \times 9 = 36\) inner
functions---matching the 36-crate ring architecture
of the trios workspace (\cite{ch:gf16-algebra},
issue \#238) \cite{kat-kolmogorov1957,kat-arnold1957}.

The discretisation proceeds as follows. Each
continuous inner function \(\psi_{p,q}\) is replaced
by a GF16 lookup table indexed by the
GoldenFloat exponent \(\hat{E}_p \in \{0, \ldots,
L_7 - 1\} = \{0, \ldots, 28\}\). The outer
functions \(\phi_q\) are realised as
Fibonacci-weighted accumulations:

\[
\phi_q(s) = \frac{1}{\sqrt{3}} \sum_{k=0}^{F_{17}-1} w_{q,k} \cdot s \cdot \varphi^{-k},
\]

where the normalisation constant
\(1/\sqrt{3} = 1/\sqrt{\varphi^2 + \varphi^{-2}}\)
is the Golden LayerNorm scale (Ch.17, Definition
3.2), ensuring that the anchor identity
\(\varphi^2 + \varphi^{-2} = 3\) is preserved at
every accumulation step.

\subsection{GF16 Inner-Function
Realisation}\label{fa_23:gf16-inner-function}

\textbf{Definition 8.2} (KAT--GF16 inner function).
For exponent index \(p \in \{1,2,3,4\}\) and
channel \(q \in \{0, \ldots, 8\}\), the discretised
inner function is

\[
\hat\psi_{p,q}(x_p) = \mathrm{GF16\_MAC}\!\bigl(\mathrm{LUT}_{p,q}[\hat{E}_p],\; \varphi^{(q \bmod 3) - 1}\bigr),
\]

where \(\mathrm{GF16\_MAC}\) is the single-cycle
multiply--accumulate operation on the QMTech
XC7A100T fabric (92 MHz, 0 DSP) and
\(\mathrm{LUT}_{p,q}\) is a 29-entry table
(addressed by \(\hat{E}_p \in \{0,\ldots,28\}\),
bounded by the Lucas sentinel \(L_7 = 29\)) whose
entries are GF16 constants in the safe domain
(INV-3, Ch.6).

\textbf{Proposition 8.3} (Inner-function LUT
budget). The total storage for all \(4 \times 9 =
36\) inner-function LUTs is \(36 \times 29 = 1044\)
GF16 entries, each 16 bits. Total BRAM requirement:
\(1044 \times 2\,\text{B} = 2088\) bytes, which fits
within a single 18 Kb BRAM tile on the XC7A100T
(the tile capacity is \(2048 \times 9 = 18432\)
bits, and \(1044 \times 16 = 16704\) bits
\(\leq 18432\)).

\subsection{Three-Band Outer-Function
Decomposition}\label{fa_23:three-band-outer}

The nine outer channels \(\phi_0, \ldots, \phi_8\)
partition into three groups of three, corresponding
to the three GF16 exponent bands induced by the
anchor identity
\(\varphi^2 + \varphi^{-2} = 3\):

\begin{enumerate}
\item \textbf{Sub-unity band} (\(q = 0, 1, 2\)):
outer scale \(\varphi^{-2} \approx 0.382\). These
channels capture the sub-normal GF16 regime where
\(\hat{E} < B\) (Ch.6, INV-3).
\item \textbf{Unity band} (\(q = 3, 4, 5\)):
outer scale \(1\). These channels operate within
the primary GF16 mantissa range.
\item \textbf{Super-unity band} (\(q = 6, 7, 8\)):
outer scale \(\varphi^{2} \approx 2.618\). These
channels handle overflow-prone accumulation
products, normalised by the
\(\varphi^{2} + \varphi^{-2} = 3\) invariant to
prevent saturation.
\end{enumerate}

\textbf{Proposition 8.4} (Band-cardinality
symmetry). Each band contains exactly three
channels. The cardinality \(3 = \varphi^2 +
\varphi^{-2}\) is the anchor identity, ensuring
that the three bands are equiproportioned in the
outer-function decomposition.

\subsection{KAT Forward-Pass Latency on the GF16
Pipeline}\label{fa_23:kat-forward-pass-latency}

The KAT--GF16 forward pass consists of two stages:
inner-function evaluation (parallel across \(p\))
followed by outer-function accumulation (sequential
across \(q\)).

\textbf{Inner stage.} Each of the four exponent
indices \(x_1, \ldots, x_4\) addresses its
nine LUTs in parallel. On the FPGA, this is a
single-cycle BRAM read at 92 MHz. Latency:
\(1 \times 10.87\,\text{ns} \approx 10.9\) ns.

\textbf{Outer stage.} The nine channels are
accumulated sequentially using the GF16 MAC unit.
Each accumulation requires one MAC cycle
(\(10.87\) ns). Total outer latency:
\(9 \times 10.87 \approx 97.8\) ns.

\textbf{Total KAT--GF16 forward-pass latency:}
\(10.9 + 97.8 = 108.7\) ns per evaluation, well
within the \(L_7 = 29\) cycle budget
(\(29 \times 10.87 = 315.2\) ns) of the
period-locked runtime monitor (Ch.24).

\subsection{Approximation Error and the Anchor
Invariant}\label{fa_23:kat-approximation-error}

The discretisation from continuous KAT inner
functions to 29-entry GF16 LUTs introduces
quantisation error bounded by the GF16 machine
epsilon in the unity band:

\[
\varepsilon_{\mathrm{KAT}} \leq \frac{9 \cdot \varepsilon_{\mathrm{GF16}}}{\sqrt{\varphi^2 + \varphi^{-2}}} = \frac{9 \cdot 2^{-10}}{\sqrt{3}} \approx \frac{9 \cdot 9.77 \times 10^{-4}}{1.732} \approx 5.08 \times 10^{-3}.
\]

The normalisation by \(\sqrt{3}\)---the square root
of the anchor identity---ensures that the error
bound is invariant under the three-band
decomposition. By the Golden LayerNorm property
(Ch.17), this bound is preserved across MCP
tool-call boundaries (Ch.23, Theorem 2.5), so the
KAT--GF16 representation composes with the agentic
inference pipeline without error amplification.

\subsection{Connection to the Ring
Architecture}\label{fa_23:kat-ring-connection}

The 36 inner functions of the KAT--GF16
decomposition correspond one-to-one to the 36 crate
rings of the trios workspace (issue \#238). Each
ring encapsulates the LUT, MAC, and normalisation
logic for one \((p, q)\) pair. The ring ordering
follows the dependency structure: rings
\(\mathrm{XX}\text{-}00\) through \(\mathrm{XX}\text{-}03\)
provide the inner-function LUTs, and rings
\(\mathrm{XX}\text{-}04\) through
\(\mathrm{XX}\text{-}08\) provide the outer-function
accumulators. The bridge ring
\(\mathrm{BR}\text{-}\mathrm{OUTPUT}\) assembles
the nine channel outputs into the final GF16 result,
applying the anchor normalisation
\(1/\sqrt{\varphi^2 + \varphi^{-2}} = 1/\sqrt{3}\).

This correspondence is not merely organisational:
it ensures that the formal verification of each ring
(AGENTS.md, invariant I5) doubles as verification
of the corresponding KAT component, and that the
\texttt{arch-guard} CI check validates the KAT
decomposition structure at every commit.

\section{References (KAT
Addendum)}\label{fa_23:references-kat}

[K14] Kolmogorov, A. N. (1957). On the
representation of continuous functions of many
variables by superposition of continuous functions
of one variable and addition.
\emph{Doklady Akademii Nauk SSSR}, 114, 953--956.

[K15] Arnold, V. I. (1957). On functions of three
variables. \emph{Doklady Akademii Nauk SSSR},
114, 679--681.

[K16] Liu, Z. et al.~(2024). KAN: Kolmogorov--Arnold
Networks. \emph{arXiv:2404.19756}.
\url{https://arxiv.org/abs/2404.19756}

[K17] This dissertation, Ch.6 --- GF16 Arithmetic
and Field Structure (INV-3 safe domain,
\(L_7 = 29\) exponent bound).

[K18] This dissertation, Ch.17 --- Ablation Matrix
(Golden LayerNorm, \(1/\sqrt{3}\) normalisation).

[K19] This dissertation, Ch.24 --- Period-Locked
Runtime Monitor (IGLA RACE scheduling, \(L_7 = 29\)
cycle budget).

[K20] gHashTag/trios, issue \#238 --- Ring
Architecture 36 Crates.

\begin{figure}[h]
  \centering
  \includegraphics[width=0.85\textwidth]{W8-23-gf16-clock-mcp.png}
  \caption{THE FIELD / THE CLOCK / THE MCP. Anchor: $\varphi^2+\varphi^{-2}=3$. [AI-generated illustration, style anchor: title-page triptych]}
  \label{fig:wave8-W8-23-gf16-clock-mcp}
\end{figure}

